\section{The local inequality needed by Ideas 5 and 6}\label{sec:local}
The next lemma is the one finite input to the strongest argument.
We first give its mathematical proof from three exact finite facts,
then justify how those facts are proved exhaustively in
Section~\ref{sec:computation}. This keeps the structure visible without
leaving a finite assertion unaccounted for.

\subsection{Endpoint slices and their closure}
For a family \(H\) on \(k\) alternatives, a bottom slice fixes the
last label \(x\) and then deletes it; a top slice fixes and deletes
the first label. The slice retains only orders having that specified
endpoint, so it is not the same as ordinary deletion from all orders.
At either end the slice sizes sum to \(|H|\).

Every nonempty endpoint slice of a member of \(\F_k\) lies in
\(\F_{k-1}\). Internal constraints remain unchanged. A triple with
the deleted endpoint and two surviving labels becomes a fixed pair
comparison, a tautology or a contradiction. Existing comparisons
involving the endpoint likewise become constant. Contradictions give
an empty slice. These observations also apply after fixing an ordered
prefix and suffix, one endpoint deletion at a time.
\lean{\path{CDCert.K6.botP}, \path{topP}, \path{slice_bot},
\path{slice_top}. The list-level \path{IsDom} represents full
extreme-law-plus-pair solution sets, with no requirement on values
at lists that are not permutations. \path{DomE} adds families empty
on the permutations. Empty slices are allowed in all sums.}

For a seven-alternative family \(J\), let \(a_j(J)\) count orders
with an enabled swap at positions \(j,j+1\), \(1\leq j\leq6\), and set
\[
 C(J)=\frac{a_1+2a_2+2a_3+2a_4+2a_5+a_6}{10}.
\]
If we sum \(A(H)\) over all top and bottom six-slices, the first
and last swaps occur once, and the four others twice. Thus
\begin{equation}\label{eq:costcount}
 \sum_{\text{top and bottom six-slices }H}A(H)=10C(J).
\end{equation}
\lean{\path{CDCert.K6.adjAt} uses swap indices \(0,\ldots,k-2\).
\path{adj} counts \(A\), and \path{cost10} is exactly \(10C\).
\path{adj6_split}, \path{cost10_split} prove the slice-count
identities. The semantic-to-list preservation of \(A\) is supplied
by the bridge in Section~\ref{sec:bridge}.}

\begin{example}[What the weighted count means]
On seven alternatives, allow either order inside each pair
\(\{1,2\},\{3,4\},\{5,6\}\), keep these three blocks in that
order, and put 7 last. There are \(2^3=8\) orders. The enabled swaps
are exactly at positions 1, 3 and 5, in every order, so
\[
 (a_1,a_2,a_3,a_4,a_5,a_6)=(8,0,8,0,8,0),\qquad
 C=\frac{8+2\cdot8+2\cdot8}{10}=4.
\]
This is a full extreme-law-plus-pair family: all cross-block comparisons
are fixed, and every triple has a compatible extreme law. Here
\(A=24\), whereas \(C=4\); they measure different things.
The weights in \(C\) are chosen so that sliding windows later count
each interior swap with total weight one.
\end{example}

\begin{theorem}[Finite seven-window estimate]\label{thm:finite}
There is a specified rational function \(q\) on six-alternative
extreme-law-plus-pair families, extended by zero to the empty family,
such that every \(J\in\F_7\), and also the empty family, satisfies
\[
 \boxed{C(J)\leq\frac58|J|+
       \sum_{\text{bottom }H}q(H)-\sum_{\text{top }H}q(H).}
\]
\end{theorem}
\lean{\path{CDCert.K6.finite_lemma} in
\path{CDCert/K6/SevenFinal.lean}. The function is fixed by
\path{qD}; it is not chosen separately for each \(J\). The theorem
is integer-scaled, as detailed below. It applies to \texttt{DomE 7 J},
including every fixed pair choice and the empty family.}

\subsection{The weights, corrections and exact facts}
It helps to view the next definitions as an accounting scheme.
The quantity \(s\) is the slack in a preliminary inequality. Splitting
it into \(\theta_b\) and \(\theta_t\) allocates that slack to the two
ends of a slice. The correction \(p\) transfers value between those
two accounts without changing their sum. The quantity \(q\) is what
remains after the internal contributions cancel. The calculations
below establish these descriptions exactly.

The certificate specifies rational weights \(\phi\) on five-label
isomorphism types and nine exceptional correction values. These are
explicit finite data, included in the companion Lean source; they
are not additional existence assumptions. Put
\[
 b(H)=\sum_{\text{bottom five-slices }P}|P|\phi(P),\qquad
 t(H)=\sum_{\text{top five-slices }P}|P|\phi(P),
\]
\[
 s(H)=\frac{12}{19}|H|+b(H)-t(H)-\frac15A(H),\qquad
 \eta=\frac1{152}.
\]
A specified correction \(p(H)\) defines
\[
 \theta_b(H)=s(H)/2+p(H),\quad
 \theta_t(H)=s(H)/2-p(H),\quad q(H)=(b(H)+t(H))/2+p(H).
\]
Outside a set \(\mathcal S\) of nine six-label types,
\(p(H)=2\eta(r_t(H)-r_b(H))\). Here \(r_b\) counts bottom
five-slices of one specified five-label type and \(r_t\) counts
top five-slices of a second specified type. On \(\mathcal S\), the
nine listed constants replace this formula. A further set
\(\mathcal T\) consists of five types, disjoint from \(\mathcal S\).

The exact finite facts are:
\begin{enumerate}[label=F\arabic*.]
\item For the eight-order type \(R_8\in\mathcal T\), \(\theta_b=\theta_t=0\).
\item Outside \(\mathcal T\cup\mathcal S\),
\[
 \theta_b(H)\geq\frac\eta2|H|+4\eta r_t(H),\qquad
 \theta_t(H)\geq\frac\eta2|H|+4\eta r_b(H).
\]
\item If some six-slice of \(J\) has type in
\((\mathcal T\setminus\{R_8\})\cup\mathcal S\), then
\[
 \Lambda(J):=\sum_{\text{bottom }H}\theta_b(H)
             +\sum_{\text{top }H}\theta_t(H)\geq\eta|J|.
\]
\end{enumerate}
The type \(R_8\in\mathcal T\) has eight orders on six alternatives:
\[
 R_8=\{524130,524310,534120,534210,
       542130,542310,543120,543210\}.
\]
The specified five-types in \(r_b,r_t\) are respectively obtained
by deleting its first and its last label. Its orders all start in 5
and end in 0, and direct calculation gives \(r_b(R_8)=r_t(R_8)=0\).
The subscript denotes its eight orders, not eight alternatives.

\lean{Definitions: \path{CDCert.K6.potW}, \path{bD}, \path{tD},
\path{sD}, \path{pD}, \path{thB}, \path{thT}, \path{qD}, \path{lamD},
\path{rb}, \path{rt}, \path{IsS}, \path{excRep6}. The two five-types
have formal indices 866 and 41. The combined six-family check is
\path{CDCert.K6.six_sem}; the regular receiver inequalities are
\path{receiver_slack}, \path{thB_bound} and \path{thT_bound}.
F3 is proved by \path{seven_bot}, \path{seven_top} after
\path{relabel_slice}, using all 26 theorems
\path{seven_fam_0} through \path{seven_fam_25}.
\path{r8_th} proves F1; \path{r8_facts} and \path{r8_slices} provide the remaining stated R8 properties.}

\subsection{Idea 10: pay for the remaining exceptional slices}
\begin{proof}[Completion of the proof of Theorem~\ref{thm:finite}]
F3 already proves \(\Lambda(J)\geq\eta|J|\) unless every six-slice
is regular, empty, or of type \(R_8\). For a regular slice, F2 supplies
a baseline \(\eta|H|/2\) and extra contributions. An \(R_8\) slice
has value zero by F1 and is missing baseline \(4\eta\).

Suppose a bottom \(R_8\) slice is obtained by fixing last label \(z\).
All its eight orders have the same first label \(w\). They therefore
also belong to one top receiver slice at \(w\). In that receiver,
fixing last label \(z\) gives the five-type counted by \(r_b\), so
F2 supplies an extra \(4\eta\) in its \(\theta_t\) value. The receiver
cannot be \(R_8\), whose \(r_b\) is zero, and cannot be empty. It is
therefore regular. Different bottom \(R_8\) slices have different
last labels, hence charge different receiver/last-label pairs. They
do not reuse a contribution. Reverse top and bottom for top \(R_8\)
slices; those charges use the separate \(\theta_b\) contributions.

All missing baselines are thus paid. At each end the slice sizes sum
to \(|J|\), giving total baseline \(\eta|J|\). Therefore
\(\Lambda(J)\geq\eta|J|\) in every case.

Now expand \(\Lambda\). The terms on five-slices obtained by deleting
one first and one last label cancel, because the two deletions
commute, with the corresponding labels canonically relabelled.
Equation~\eqref{eq:costcount} gives the remaining swap contribution.
The resulting exact identity is
\[
 \Lambda(J)=\frac{12}{19}|J|-C(J)+
       \sum_{\text{bottom }H}q(H)-\sum_{\text{top }H}q(H).
\]
Finally \(12/19-\eta=(96-1)/152=5/8\), proving the claimed inequality.
The empty-family case has every term zero.
\end{proof}
\lean{\path{CDCert.K6.lam_receiver} is the receiver argument;
\path{seven_finite} combines it with F3. The cancellation identity
is \path{lamD_eq}. Thus neither the exceptional R8 family nor
coverage of all seven-families is left as a premise of
\path{finite_lemma}. Section~\ref{sec:computation} supplies the
finite proofs of F1--F3, completing the dependency chain.}

\subsection{The exact integer normalisation}\label{sec:scaling}
Let \(d=8\,846\,117\,280\). Formal variables are normalised as follows:
\begin{center}
\begin{tabular}{@{}ll@{}}\toprule
Formal quantity & Mathematical quantity\\\midrule
\texttt{cD5} & \(d/5=1\,769\,223\,456\)\\
\texttt{cK}, \texttt{cE} & \(12d/19=5\,587\,021\,440\), \(d/152=58\,198\,140\)\\
\texttt{cost10} & \(10C\)\\
\texttt{sD}, \texttt{bD}, \texttt{tD} & \(ds,db,dt\)\\
\texttt{pD}, \texttt{thB}, \texttt{thT}, \texttt{qD} & \(2dp,2d\theta_b,2d\theta_t,2dq\)\\
\texttt{lamD} & \(2d\Lambda\)\\\bottomrule
\end{tabular}
\end{center}
Multiplying Theorem~\ref{thm:finite} by \(2d\) gives literally
\[
 (d/5)\,\texttt{cost10}(J)\leq
 11\,057\,646\,600\,|J|+
 \sum_{\text{bottom }H}\texttt{qD}(H)-
 \sum_{\text{top }H}\texttt{qD}(H).
\]
This is the statement of \path{CDCert.K6.finite_lemma}. Counts are
cast into integers before subtraction. Consequently there is no
truncated subtraction of a possibly negative potential in this step.
